% ============================================================
%  chapters/ch06-evaluation-dynamics.tex
% ============================================================

\chapter{Evaluation Dynamics}

\chapterepigraph{%
  Evaluation is not search.\\
  It is descent along a gradient\\
  that was always already there.
}{Flyxion, \textit{Semantic Infrastructure}}

\sidenote{App.\,C §C.6--C.10\\for formal proofs\\of termination\\and confluence}

Chapter~5 defined the static syntax of Spherepop. This chapter defines
the dynamics: how expressions change over time as they are evaluated.

\section{The Pop Operation}

The atomic unit of dynamics is the pop.

\begin{definition}{Redex}{redex}
  A sphere $s = \sphere{L, [c_1,\ldots,c_n]}$ in expression $E$ is
  a \emph{redex} (reducible expression) if:
  \begin{enumerate}
    \item All contents $c_i$ are atomic (no content sphere contains
          further spheres), and
    \item The operator $L$ is defined on the values of $c_1,\ldots,c_n$.
  \end{enumerate}
  A redex is the innermost sphere ready to pop.
\end{definition}

\begin{definition}{Pop Step}{pop-step}
  A \emph{pop step} replaces a redex
  $s = \sphere{L, [c_1,\ldots,c_n]}$ with the atomic sphere
  $\sphere{L(c_1,\ldots,c_n)}$, where $L(c_1,\ldots,c_n)$ is the
  result of applying operator $L$ to the atomic values.
  We write $E \to_\text{pop} E'$ when $E'$ is obtained from $E$
  by one pop step.
\end{definition}

\section{Refusal}

Not every innermost sphere can pop. Refusal is suspended evaluation.

\begin{definition}{Refusal}{refusal}
  A sphere $s$ \emph{refuses} when:
  \begin{enumerate}
    \item \textbf{Variable refusal:} $s$ contains a free variable
          (unbound symbol). $s$ cannot pop until the variable is bound.
    \item \textbf{Domain refusal:} $s$'s operator is undefined on
          its current arguments (e.g., division by zero, $\sqrt{-1}$
          in real arithmetic). $s$ cannot pop until the arguments
          are modified.
    \item \textbf{Protection:} $s$ carries a protection marker $\protect{}$
          that prevents popping regardless of content.
  \end{enumerate}
  A refusing sphere is written $\refuse{}_s$, with subscript indicating
  the reason.
\end{definition}

\begin{remark}
  Refusal models lazy evaluation, closures, suspended computations, and
  blocked threads. A refusing sphere is not an error — it is a sphere
  waiting for its preconditions to be satisfied. Many functional
  programming patterns (promises, futures, generators) are formalized
  by refusal and the subsequent binding that resolves it.
\end{remark}

\section{Binding}

\begin{definition}{Binding Step}{binding}
  A \emph{binding step} assigns a value $v$ to a free variable $x$
  throughout an expression $E$, written $E[x := v]$. After binding,
  any sphere that was refusing due to $x$ may become a redex.
\end{definition}

\begin{example}
  The sphere $\sphere{+, [x, \sphere{3}]}$ refuses because $x$ is free.
  After binding $x := 5$, it becomes $\sphere{+, [\sphere{5}, \sphere{3}]}$,
  which is a redex that pops to $\sphere{8}$.
\end{example}

\section{Collapse Operators}

\begin{remark}
  \textbf{The Spherepop operators in quantum context.} The four operators
  Pop, Bind, Collapse, and Refuse acquire a striking interpretation when
  applied to quantum measurement (the HYDRA/Spherepop connection):

  \begin{itemize}
    \item \textbf{Pop (Create):} an irreversible quantum event — the
          moment a photon strikes a detector, puncturing the continuous
          wavefunction with a discrete, permanent fact.
    \item \textbf{Bind (Join):} correlation of two compatible quantum
          events into a joint structure. The beginning of quantum geometry.
    \item \textbf{Collapse (Identify):} grouping different causal
          histories that lead to the same macroscopic outcome — the
          many paths that produce the same measurement result.
    \item \textbf{Refuse (Filter):} the topological filter that the RSVP
          entropy field applies: any cluster of events that exceeds
          the available trajectory volume, or that contradicts the
          vector flow field, is simply refused. This is the dynamical
          mechanism that replaces the Born rule — rather than an
          appended postulate, measurement collapse follows from the
          geometry of the accessibility landscape refusing inconsistent
          branches.
  \end{itemize}

  In this reading, Schrödinger's equation is not the foundation of
  quantum mechanics but its continuum limit: the smooth approximation
  valid between irreversible events. The events (Pops) are primary;
  the wavefunction is their residue.
\end{remark}


Beyond individual pops, Spherepop supports compound operations that
collapse multiple levels simultaneously.

\begin{definition}{Collapse Operator}{collapse-op}
  A \emph{collapse operator} $\collapse{k}$ applied to an expression $E$
  performs $k$ pop steps simultaneously, collapsing the $k$ deepest
  levels. Formally:
  \[
    \collapse{k}(E) = E' \text{ where } E \to_\text{pop}^k E'.
  \]
\end{definition}

Collapse operators are useful for describing batch evaluation — the
behavior of a compiler that reduces an entire expression in one pass
rather than step by step.

\begin{definition}{Full Collapse}{full-collapse}
  The \emph{full collapse} $\collapse{\infty}(E)$ is the unique terminal
  expression reachable from $E$ by exhaustive popping (if it exists).
\end{definition}

\section{Reduction Sequences}

\begin{definition}{Reduction Sequence}{reduction-seq}
  A \emph{reduction sequence} from $E$ is a (possibly infinite) sequence
  \[
    E = E_0 \to_\text{pop} E_1 \to_\text{pop} E_2 \to_\text{pop} \cdots
  \]
  of pop steps. A reduction sequence \emph{terminates} at $E_n$ if
  $E_n$ contains no redexes.
\end{definition}

\begin{theorem}{Termination for Arithmetic}{termination-arith}
  Every finite, ground (variable-free), well-typed arithmetic Spherepop
  expression has a terminating reduction sequence.
\end{theorem}

\begin{proof}
  By strong induction on the depth $d(E)$. If $d(E) = 0$, $E$ is atomic
  and already terminal. If $d(E) \geq 1$, some innermost sphere $s$ is
  a redex (since $E$ is ground and well-typed). Popping $s$ strictly
  reduces the number of non-atomic spheres. Since this count is
  non-negative and integer-valued, the sequence terminates.
\end{proof}

\section{Termination Conditions}

The termination theorem above requires three conditions that are worth
examining:

\begin{itemize}
  \item \textbf{Finite:} infinite Spherepop diagrams (co-recursive
        expressions) may not terminate. Streams, infinite lists, and
        co-inductive data types require lazy evaluation and controlled
        refusal.
  \item \textbf{Ground:} free variables cause refusal. A variable-free
        expression has no refusing spheres (assuming well-typedness)
        and must terminate.
  \item \textbf{Well-typed:} domain refusal (division by zero, etc.)
        can prevent termination even in ground expressions.
        Type safety rules out domain errors statically.
\end{itemize}

The interaction of these three conditions defines the boundary between
\emph{total} computation (always terminates) and \emph{partial} computation
(may not terminate). The Halting Problem is precisely the question of
which Spherepop expressions have terminating reduction sequences — a
question that is undecidable in general.

\section{Confluence}

A key property of Spherepop evaluation is that the choice of which
redex to pop first does not affect the final result.

\begin{theorem}{Confluence}{confluence-ch6}
  For any Spherepop expression $E$ over a confluent operator set
  (one where all operators are deterministic and total on their domain),
  any two reduction sequences from $E$ reach the same terminal expression.
\end{theorem}

\begin{proof}[Sketch]
  Any two redexes in $E$ are either nested (one is inside the other)
  or disjoint (neither contains the other). If nested, the outer redex
  depends on the inner, so they cannot both be redexes simultaneously —
  the outer will only become a redex after the inner pops. If disjoint,
  popping either one does not affect the other, so their pops commute.
  By the Church–Rosser property for disjoint reductions, all sequences
  reach the same terminal.
\end{proof}

Confluence means evaluation is \emph{deterministic in result} even when
it is \emph{non-deterministic in strategy}. Different evaluation orders
may take different numbers of steps, but they always arrive at the same place.

\begin{exercise}
  Give an example of a Spherepop expression with two distinct redexes
  that can be popped in either order. Show that both orders reach the
  same terminal expression, verifying confluence in this case.
\end{exercise}

\begin{exercise}
  Define a Spherepop expression that refuses due to each of the three
  reasons (variable, domain, protection). For each, describe what
  additional step would resolve the refusal and allow evaluation
  to continue.
\end{exercise}

\begin{exercise}[title={Harder}]
  Construct a Spherepop expression that does not terminate. (Hint:
  consider an operator that expands its arguments rather than reducing
  them — analogous to the $\Omega$ combinator in the $\lambda$-calculus.)
  What happens to the depth of the expression as evaluation proceeds?
\end{exercise}
